% ---
% id: foundations/def:tiers
% kind: def
% title: "The tiers of a reported key rate"
% summary: >-
%   A reported key rate carries one of three tiers: estimate, a returned
%   point's objective, claiming no bound; first-order, a lower bound on the
%   minimum under named hypotheses; or certified, such a bound with every step
%   from the data exact or rigorously enclosed, no hypothesis about arithmetic.
% prominence: standard
% status: implemented
% uses: []
% mentions: [meat-core/lem:estimate, meat-core/frank-wolfe/thm:firstorder, meat-core/frank-wolfe/lem:certified-bracket]
% assumes: []
% refs: []
% provenance: newly-derived
% implements: certiqs_key.foundations.assurance:Assurance
% lean:
% ---

In this repository the estimate is the objective value of
\xref{meat-core/lem:estimate}, and the first-order number is the linearisation
bound of the Frank--Wolfe solver, \xref{meat-core/frank-wolfe/thm:firstorder},
whose hypotheses include lower bounds on the two linear minima it reads. A
run evaluates that bound in binary64, so a first-order number rests, beside
the hypotheses of the theorem, on arithmetic that no claim bounds. The
certified number is the same bound formed again at a run's final points from
enclosures of the gradient
(\xref{meat-core/frank-wolfe/lem:certified-bracket}), every step from the
declaration's data exact or carried out in an arithmetic that bounds its own
rounding, so that nothing is assumed of how it was computed.

\begin{definition}\label{def:tiers}
A \term[one of three labels, estimate, first-order or certified, that a
reported key-rate number carries]{tier} of a reported key-rate number is one
of the following three labels.
\begin{enumerate}[label=(\roman*),leftmargin=2.4em]
\item \term[the tier of a key-rate number that is the value of the objective
  at the point an optimiser returned, claimed as no bound]{estimate}: the
  number is the value of the objective at the point
  an optimiser returned. It is reported as an estimate of the minimum of the
  objective and is claimed as no bound.
\item \term[the tier of a key-rate number that is a lower bound on the minimum
  of the objective, established by a claim under hypotheses that claim
  names]{first-order}: the number is a lower bound on the minimum of the
  objective, established by a claim under hypotheses that claim names.
\item \term[the tier of a lower bound on the minimum of the objective of a
  declaration, formed with every arithmetic step from that declaration's
  data exact or rigorously enclosed, so that none of its hypotheses
  concerns arithmetic]{certified}: the number is a lower bound on the minimum
  of the objective of the declaration it is formed on, established by a claim
  under hypotheses that claim names, and every arithmetic step between that
  declaration's data and the number is exact or
  \term[said of a computation each of whose operations returns a set of
  numbers, such as an interval or a ball, containing the exact result of the
  operation at every value its inputs' sets contain]{rigorously enclosed}: a
  computation is rigorously enclosed when each of its operations returns a
  set of numbers, such as an interval or a ball, that contains the exact
  result of the operation at every value the sets of its inputs contain, so
  that the set it ends with contains the exact value, at the exact data, of
  the quantity it computes. None of those hypotheses is an assumption about
  how a number was computed, that is, none is an assumption of the category
  numerical-trust.
\end{enumerate}
The tiers are ordered, estimate below first-order below certified, and a
tier is \emph{stronger} than the tiers below it. A certified number is a
first-order number as well. There is no other tier.
\end{definition}
